% =====================================================================
%  MBLP131 -- İşletim Sistemleri Temelleri
%  Hafta 15 Ders Notları: Final Sınavı Örnek Kâğıdı
%  XeLaTeX ile iki kez derleyin:  xelatex MBLP131_Hafta15_Notlar.tex
%  ./figures/ ve ../code/ klasörlerine ihtiyaç duyar.
% =====================================================================
\documentclass[11pt,a4paper]{article}
\newcommand{\weeknum}{15}
% ---------- Page & fonts ---------------------------------------------
\usepackage[a4paper,margin=2.3cm,headheight=15pt]{geometry}
\usepackage{fontspec}
\usepackage{polyglossia}
\setdefaultlanguage{turkish}
\IfFontExistsTF{Calibri}{\setmainfont{Calibri}}{\setmainfont{Carlito}}
\IfFontExistsTF{Cambria}{\newfontfamily\headfont{Cambria}}{\newfontfamily\headfont{Caladea}}
\IfFontExistsTF{Consolas}{\setmonofont[Scale=0.95]{Consolas}}{\setmonofont[Scale=0.88]{DejaVu Sans Mono}}
\usepackage{amsmath,amssymb}
\usepackage{microtype}
\linespread{1.12}
\setlength{\parskip}{0.55em}
\setlength{\parindent}{0pt}

% ---------- Colours (course design system) ----------------------------
\usepackage{xcolor}
\definecolor{dark}{HTML}{212121}
\definecolor{charcoal}{HTML}{36454F}
\definecolor{slate}{HTML}{50808E}
\definecolor{mint}{HTML}{02C39A}
\definecolor{mintdeep}{HTML}{028090}
\definecolor{light}{HTML}{F2F2F2}
\definecolor{muted}{HTML}{6B7679}

% ---------- Packages ---------------------------------------------------
\usepackage{graphicx}
\usepackage{float}
\usepackage{needspace}
\graphicspath{{figures/}}
\usepackage[font=small,labelfont={bf,color=mintdeep},justification=centering]{caption}
\usepackage{booktabs,longtable,array,tabularx}
\usepackage{enumitem}
\setlist{itemsep=0.25em,topsep=0.3em}
\usepackage{tikz}
\usetikzlibrary{arrows.meta,positioning,calc,shapes.geometric}
\usepackage{pgfplots}
\pgfplotsset{compat=1.17}
\usepackage[most]{tcolorbox}
\usepackage{titlesec}
\usepackage{fancyhdr}
\usepackage{xurl}
\usepackage[colorlinks=true,linkcolor=mintdeep,urlcolor=mintdeep,citecolor=mintdeep]{hyperref}

% ---------- Headings ----------------------------------------------------
\titleformat{\section}{\headfont\Large\bfseries\color{charcoal}}{\color{mint}\thesection}{0.7em}{}
\titleformat{\subsection}{\headfont\large\bfseries\color{charcoal}}{\color{mintdeep}\thesubsection}{0.6em}{}
\titleformat{\subsubsection}{\headfont\normalsize\bfseries\color{mintdeep}}{}{0em}{}
\titlespacing*{\section}{0pt}{1.6em}{0.6em}
\titlespacing*{\subsection}{0pt}{1.2em}{0.4em}

% ---------- Header / footer -----------------------------------------------
\pagestyle{fancy}
\fancyhf{}
\fancyhead[L]{\small\color{muted}MBLP131 --- İşletim Sistemleri Temelleri}
\fancyhead[R]{\small\color{muted}Hafta \weeknum{} Ders Notları}
\fancyfoot[C]{\small\color{muted}\thepage}
\renewcommand{\headrulewidth}{0pt}

% ---------- Boxes ---------------------------------------------------------
\newtcolorbox{keyidea}[1][Key idea]{enhanced,breakable,colback=charcoal,colframe=charcoal,
  coltext=white,fonttitle=\headfont\bfseries\color{mint},title={#1},
  attach title to upper={\par\smallskip},boxrule=0pt,arc=3pt,left=10pt,right=10pt,top=8pt,bottom=8pt}
\newtcolorbox{example}[1][Worked example]{enhanced,colback=light,colframe=light,
  fonttitle=\headfont\bfseries\color{mintdeep},title={#1},attach title to upper={\par\smallskip},
  boxrule=0pt,arc=3pt,left=10pt,right=10pt,top=8pt,bottom=8pt}
\newtcolorbox{tryit}[1][Try it yourself]{enhanced,breakable,colback=white,colframe=mint,
  fonttitle=\headfont\bfseries\color{mintdeep},title={#1},attach title to upper={\par\smallskip},
  boxrule=0.8pt,arc=3pt,left=10pt,right=10pt,top=8pt,bottom=8pt}
\newtcolorbox{histnote}[1][Historical note]{enhanced,breakable,colback=white,colframe=slate!40,
  fonttitle=\headfont\bfseries\color{slate},title={#1},attach title to upper={\par\smallskip},
  boxrule=0.5pt,arc=3pt,left=10pt,right=10pt,top=8pt,bottom=8pt,fontupper=\small}

% Photo helper: fits image inside a box, never distorted (keepaspectratio)
\newcommand{\photo}[3]{\includegraphics[width=#1,height=#2,keepaspectratio]{#3}}

\newcommand{\term}[1]{\textbf{\color{mintdeep}#1}}
\newcommand{\cmd}[1]{\texttt{#1}}


% ---------- Code listings ----------------------------------------------------
\usepackage{listings}
\lstdefinestyle{course}{
  basicstyle=\ttfamily\small, columns=fullflexible, keepspaces=true,
  backgroundcolor=\color{dark}, rulecolor=\color{dark}, frame=single, framesep=6pt,
  xleftmargin=8pt, xrightmargin=8pt, aboveskip=0.8em, belowskip=0.6em,
  basicstyle=\ttfamily\small\color{light},
  keywordstyle=\color{mint}\bfseries, commentstyle=\color{white!60}\itshape,
  stringstyle=\color{orange!70!yellow}, numberstyle=\tiny\color{muted},
  numbers=left, numbersep=10pt, showstringspaces=false, breaklines=true,
  tabsize=4, upquote=true}
\lstdefinestyle{term}{style=course, numbers=none, keywordstyle=\color{light},
  commentstyle=\color{white!60}, stringstyle=\color{light}}
\lstset{style=course,language=C}
\usepackage{textcomp}

% listings 1.9 + XeLaTeX: Turkish letters above U+00FF (ı İ ğ Ğ ş Ş) must be declared as letters
\makeatletter
\lst@CCPut \lst@ProcessLetter {"130}{"131}{"11E}{"11F}{"15E}{"15F}\z@
\lst@AddToHook{SelectCharTable}{\catcode"130=13 \catcode"131=13 \catcode"11E=13
  \catcode"11F=13 \catcode"15E=13 \catcode"15F=13 }
\makeatother

\begin{document}
% ===================== MBLP131 Final Sınavı -- ÖRNEK KÂĞIT (öğrenci sürümü) =====================
% Dersin önsözüyle derlenir:  python3 ../../lib/assemble_notes.py tr 15 final_ornek_body.tex MBLP131_Final_Ornek.tex "Final Sınavı Örnek Kâğıdı"
% Her yanıt ../code içinde doğrulandı (theory_check.py, answers_sample.sh); bkz. MBLP131_Final_CevapAnahtari.
\setlength{\emergencystretch}{3em}\sloppy
\fancyhead[R]{\small\color{muted}Final Sınavı --- Örnek Kâğıt}
\setlength{\parskip}{0.35em}
\newcommand{\pts}[1]{\hfill\textcolor{mintdeep}{\textbf{[#1~puan]}}}
\newcommand{\ans}[1]{\par\smallskip\noindent\begin{tikzpicture}\draw[muted!60,rounded corners=2pt] (0,0) rectangle (\linewidth,#1);\end{tikzpicture}\par\medskip}
% Gantt şeması için boş zaman ekseni, 0..#1 ms
\newcommand{\ganttgrid}[2]{\par\noindent\begin{tikzpicture}[x=0.7cm]
\node[anchor=east,font=\small\bfseries] at (-0.2,0.31) {#2};
\draw[muted!70] (0,0) rectangle (#1,0.62);
\foreach \t in {0,...,#1} {\draw[muted!70] (\t,0) -- (\t,-0.1) node[below,font=\scriptsize,text=muted,inner sep=1pt] {\t};}
\end{tikzpicture}\par\smallskip}
% boş çerçeve tablosu: 12 başvuru, 3 çerçeve
\newcommand{\frametable}[1]{\begin{tabular}{@{}l|*{12}{>{\centering\arraybackslash}p{0.6cm}|}@{}}
\multicolumn{13}{@{}l}{\textbf{#1}}\\ \hline
başv. & 5 & 4 & 1 & 5 & 0 & 5 & 4 & 0 & 2 & 1 & 0 & 1 \\ \hline
çr.\ 0 &&&&&&&&&&&& \\ \hline
çr.\ 1 &&&&&&&&&&&& \\ \hline
çr.\ 2 &&&&&&&&&&&& \\ \hline
hata &&&&&&&&&&&& \\ \hline
\end{tabular}}

\thispagestyle{empty}
\begin{tcolorbox}[enhanced,colback=dark,colframe=dark,arc=0pt,boxrule=0pt,
  left=18pt,right=18pt,top=14pt,bottom=14pt,coltext=white]
{\small\color{mint}\textsc{MBLP131 \,·\, İşletim Sistemleri Temelleri}}\par\smallskip
{\headfont\fontsize{22}{26}\selectfont\bfseries Final Sınavı (Yarıyıl Sonu) --- Örnek Kâğıt}\par\smallskip
{\small İstanbul Okan Üniversitesi \,·\, Meslek Yüksekokulu \,·\, Elektronik Haberleşme Teknolojisi Programı}\par
{\small Öğretim üyesi: Dr.\ Öğr.\ Üyesi Hasan Oğuz \quad Tarih: [duyurulacak] \quad Süre: [90] dakika}
\end{tcolorbox}

\begin{center}
\renewcommand{\arraystretch}{1.6}
\begin{tabularx}{\linewidth}{@{}lX lX@{}}
\textbf{Ad, soyad:} & \hrulefill & \textbf{Öğrenci no:} & \hrulefill \\
\textbf{İmza:} & \hrulefill & \textbf{Sıra / salon:} & \hrulefill \\
\end{tabularx}

\vspace{0.4em}
\renewcommand{\arraystretch}{1.35}
\begin{tabular}{|l|c|c|c|c|c|c|c|}
\hline
\textbf{Bölüm} & A & B & C & D & E & F & \textbf{Toplam} \\ \hline
\textbf{Puan} & 20 & 12 & 36 & 14 & 9 & 9 & \textbf{100} \\ \hline
\textbf{Alınan} & & & & & & & \\[0.4em] \hline
\end{tabular}
\end{center}

\textbf{Kurallar.} Kitap kapalı [kesinleşecek]. [Basit, programlanamayan hesap makinesi serbesttir --- kesinleşecek.] Telefonlar kapalı. Yanıtlarınızı kutulara yazın; her hesabın tüm adımlarını gösterin (formül $\to$ sayılar $\to$ birimiyle sonuç). Dersin zamanlama uzlaşımları: bağlam değişimi süresi $0$; eşitlikler önce varış zamanına, sonra ada göre bozulur; küçük öncelik numarası yüksek öncelik demektir; Round Robin'de bir zaman diliminin dolduğu anda gelen süreç kuyruğa kesilen süreçten önce girer. TLB: $\text{EAT} = t_{\text{TLB}} + (2-\alpha)\,t_m$ formülünü kullanın. Linux soruları \texttt{\textasciitilde/lab} alıştırma klasörüne (\texttt{mcmp117\_lab\_setup.sh} ile oluşturulur) göredir; aksi belirtilmedikçe çalışma dizininiz \texttt{\textasciitilde/lab}'dır.

% ---------------------------------------------------------------------
\section*{Bölüm A --- Çoktan seçmeli \hfill (10 $\times$ 2 = 20 puan)}
\emph{Bir} yanıtı daire içine alın.
\begin{enumerate}[leftmargin=1.8em,label=A\arabic*.,itemsep=0.35em]
  \item UNIX, Bell Labs'ta hangi yıl geliştirildi? \quad (a) 1956 \quad (b) 1969 \quad (c) 1981 \quad (d) 1991
  \item Hangi işletim sistemi yapısında çekirdek modunda yalnızca küçük bir çekirdek çalışır, hizmetlerin çoğu kullanıcı alanındadır? \quad (a) monolitik \quad (b) basit \quad (c) mikroçekirdek \quad (d) modüler
  \item Hangisi süreç denetim bloğunda (PCB) \emph{saklanmaz}? \quad (a) süreç durumu \quad (b) program sayacı \quad (c) açık dosyaların listesi \quad (d) programın kaynak kodu
  \item Bir sürecin varışından işlemciyi ilk aldığı ana kadar geçen süre onun \quad (a) dönüş süresi \quad (b) bekleme süresi \quad (c) yanıt süresi \quad (d) patlama süresidir
  \item Hangi semafor işlemi çağıran iş parçacığını bloklayabilir? \quad (a) \texttt{signal(S)} \quad (b) \texttt{wait(S)} \quad (c) ikisi de \quad (d) hiçbiri
  \item Bir program çalışırken işlemcinin ürettiği adresler \quad (a) fiziksel \quad (b) mantıksal \quad (c) disk blok numaraları \quad (d) düğüm (inode) numaralarıdır
  \item Boşa çalışma (thrashing) ne demektir? \quad (a) çok yüksek işlemci kullanımı \quad (b) sürekli sayfa hataları ve düşük işlemci kullanımı \quad (c) \%100 TLB isabet oranı \quad (d) dolu bir disk
  \item \texttt{chmod 754 f} sonrasında \texttt{ls -l f} ne gösterir? \quad (a) \texttt{-rwxr-xr--} \quad (b) \texttt{-rwxr-x---} \quad (c) \texttt{-rwxrw-r--} \quad (d) \texttt{-rw-r-xr--}
  \item \texttt{cmd 2> /dev/null} \quad (a) normal çıktıyı atar \quad (b) hata iletilerini atar \quad (c) hataları \texttt{2} adlı bir dosyaya yazar \quad (d) çıktıyı sona ekler
  \item Hangi komut geçerli kabuğun arka plan işlerini listeler? \quad (a) \texttt{ps} \quad (b) \texttt{top} \quad (c) \texttt{jobs} \quad (d) \texttt{bg}
\end{enumerate}

% ---------------------------------------------------------------------
\section*{Bölüm B --- Kısa yanıt \hfill (3 $\times$ 4 = 12 puan)}
\begin{enumerate}[leftmargin=1.8em,label=B\arabic*.]
  \item Bir program \texttt{write()} çağırıyor. Çağrının çekirdeğe nasıl ulaştığını ve nasıl geri döndüğünü adım adım anlatın (API/sarmalayıcı, sistem çağrısı numarası, tuzak, mod biti, sistem çağrısı tablosu). \pts{4}
  \ans{3.2cm}
  \item Yarış durumu nedir? Örnek olarak \texttt{counter++}'ı kullanın ve kritik bölge problemine bir çözümün sağlaması gereken üç koşulu yazın. \pts{4}
  \ans{3.2cm}
  \item Boşa çalışma (thrashing) nedir ve çalışma kümesi modeli işletim sisteminin bunu önlemesine nasıl yardım eder? \pts{4}
  \ans{3.2cm}
\end{enumerate}

% ---------------------------------------------------------------------
\newpage
\section*{Bölüm C --- Sayısal problemler \hfill (36 puan)}
\begin{enumerate}[leftmargin=1.8em,label=C\arabic*.]
  \item Dört süreç (süreler ms):
  \begin{center}
  \begin{tabular}{lcccc}
  \toprule
  Süreç & P1 & P2 & P3 & P4 \\ \midrule
  Varış & 0 & 2 & 3 & 5 \\
  Patlama & 6 & 2 & 7 & 3 \\
  Öncelik & 3 & 1 & 4 & 2 \\ \bottomrule
  \end{tabular}
  \end{center}
  Her algoritma için Gantt şemasını çizin, ortalama bekleme süresini ve ortalama dönüş süresini hesaplayın: (a) FCFS, (b) kesintili Öncelik, (c) $q = 2$ ile Round Robin. \pts{12}
  \ganttgrid{18}{FCFS}
  \ganttgrid{18}{Öncelik}
  \ganttgrid{18}{RR, $q=2$}
  \ans{4.6cm}
  \item Sayfalar $2$~KB. Sayfa tablosu: sayfa $0 \to$ çerçeve $3$, $1 \to 8$, $2 \to 6$, $3 \to 1$ (dört girdi).
  (a) $5000$ mantıksal adresini fiziksel adrese çevirin. \pts{3}\\
  (b) $9000$ mantıksal adresi için ne olur? \pts{1}\\
  (c) $t_{\text{TLB}} = 10$~ns, $t_m = 90$~ns, $\alpha = 0.8$: etkin erişim süresini hesaplayın. \pts{2}\\
  (d) Boş delikler (bu sırayla) $200$, $450$, $100$, $300$, $500$~KB. İlk uyan, en iyi uyan ve en kötü uyan ile $280$~KB'lık bir isteği hangi delik alır? \pts{2}
  \ans{5.2cm}
  \Needspace{20\baselineskip}
  \item $3$ çerçeveyle $5, 4, 1, 5, 0, 5, 4, 0, 2, 1, 0, 1$ başvuru dizisi. FIFO, OPT ve LRU için çerçeve tablolarını doldurun ve sayfa hatası sayısını verin. \pts{8}
  \begin{center}
  \small\renewcommand{\arraystretch}{1.2}
  \frametable{FIFO}\quad hata: \underline{\hspace{1.2cm}}\\[0.5em]
  \frametable{OPT}\quad hata: \underline{\hspace{1.2cm}}\\[0.5em]
  \frametable{LRU}\quad hata: \underline{\hspace{1.2cm}}
  \end{center}
  \item (a) $t_m = 100$~ns, $t_{\text{fault}} = 5$~ms, sayfa hatası oranı $p = 0.0005$: etkin erişim süresini hesaplayın. RAM'den kaç kat yavaştır? \pts{3}\\
  (b) Bir dosyanın kayıtları $50$ bayt, blokları $512$ bayttır; ikisi de $0$'dan numaralanır. Kayıt $100$ hangi blokta ve hangi konumda başlar? Kayıt o bloğa sığar mı? \pts{3}\\
  (c) İşler zamanın yarısında G/Ç bekliyor ($p = 0.5$): $U(4)$'ü hesaplayın. Bir programın \%25'i seridir: Amdahl yasası $3$ çekirdekte ne kadar hızlanmaya izin verir, çok sayıda çekirdek için sınır nedir? \pts{2}
  \ans{6.0cm}
\end{enumerate}

% ---------------------------------------------------------------------
\newpage
\section*{Bölüm D --- Komutu yazın \hfill (7 $\times$ 2 = 14 puan)}
Her görev için \emph{tek} bir komut satırı yazın.
\begin{enumerate}[leftmargin=1.8em,label=D\arabic*.]
  \item \texttt{docs} dizininin gizli dosyalar dahil, okunabilir boyutlarla uzun listesi. \ans{0.9cm}
  \item \texttt{docs/notes\_01.txt} dosyasının adını \texttt{docs/meeting\_01.txt} olarak değiştirin. \ans{0.9cm}
  \item \texttt{report1.txt} \dots\ \texttt{report5.txt} boş dosyalarını kısa tek bir komutla oluşturun. \ans{0.9cm}
  \item Çalışma dizininde \texttt{logs/access.log} dosyasını gösteren \texttt{latest.log} adlı bir sembolik bağlantı oluşturun. \ans{0.9cm}
  \item \texttt{scripts/backup.sh} için (yalnızca) sahibine çalıştırma izni ekleyin; sembolik gösterimle. \ans{0.9cm}
  \item PID'i $4321$ olan süreç SIGTERM'i yok sayıyor. Onu yine de durdurun. \ans{0.9cm}
  \item \texttt{logs} dizininin ne kadar disk alanı kullandığını okunabilir biçimde gösterin. \ans{0.9cm}
\end{enumerate}

\section*{Bölüm E --- Boru hattını yazın \hfill (3 $\times$ 3 = 9 puan)}
\begin{enumerate}[leftmargin=1.8em,label=E\arabic*.]
  \item \texttt{data/cities.txt} dosyasında her satırda bir şehir var. En sık geçen şehri sayısıyla birlikte yazdırın. \ans{1.1cm}
  \item \texttt{logs/access.log} içindeki isteklerin kaçı \texttt{/login} sayfası içindi? (Bir satır şöyle görünür:\\ \texttt{10.0.1.10 - - [12/Oct/2026:08:00:00 +0300] "GET /about.html HTTP/1.1" 404 442}.) \ans{1.1cm}
  \item \texttt{data/students.csv} dosyasının başlığı \texttt{id,name,department,midterm,final}'dır. Final notu en yüksek üç öğrencinin adını ve final notunu (\texttt{name,final}) en iyiden başlayarak yazdırın. \ans{1.1cm}
\end{enumerate}

\section*{Bölüm F --- Bir betik yazın \hfill (9 puan)}
\texttt{passfail.sh} betiğini yazın. Betik tam olarak bir argüman, bir eşik değeri alır. Tam olarak bir argüman verilmezse standart hataya bir kullanım iletisi yazar ve $1$ durumuyla çıkar. Aksi hâlde \texttt{data/students.csv} içinde final notu eşiğe eşit ya da eşikten büyük olan (``passed'') ve eşikten küçük olan (``failed'') öğrencileri sayar ve örneğin \texttt{threshold 50: 30 passed, 10 failed} biçiminde yazdırır.
\ans{9.0cm}

\begin{center}\itshape --- Kâğıdın sonu. Yanıtlarınızı kontrol edin. Başarılar! ---\end{center}

\end{document}
